
\usepackage[utf8]{inputenc}
\DeclareUnicodeCharacter{03A9}{\ensuremath{\Omega}}
\usepackage[T1]{fontenc}
\usepackage{lmodern}
\usepackage{geometry}
\geometry{margin=1.2in}
\usepackage{amsmath, amsthm, amssymb, mathtools, extarrows}
\usepackage{enumitem}
\usepackage{tikz-cd}
\usepackage{quiver}
\usepackage[colorlinks=true, linkcolor=blue, citecolor=red, urlcolor=blue]{hyperref}
\bibliographystyle{abbrvurl-doi}
\usepackage{xurl}
\usepackage{xcolor}
\newcommand{\todo}[1]{\textcolor{red}{[\textbf{TODO:} #1]}}
\usepackage{soul}
\usepackage{framed}
\colorlet{cutyellow}{yellow!60}
\sethlcolor{cutyellow}
\newcommand{\cut}[1]{\hl{#1}}
\newenvironment{cutblock}{\def\FrameCommand{\fboxsep=4pt\colorbox{cutyellow}}\MakeFramed{\advance\hsize-\width\FrameRestore}}{\endMakeFramed}
\usepackage{comment}
\colorlet{proposalblue}{blue!7}
\newif\ifshowproposals
\showproposalstrue
\ifshowproposals
\specialcomment{proposal}{\def\FrameCommand{\fboxsep=6pt\colorbox{proposalblue}}\MakeFramed{\advance\hsize-\width\FrameRestore}\noindent\textbf{Proposal.}\ }{\endMakeFramed}
\else
\excludecomment{proposal}
\fi

\theoremstyle{plain}
\newtheorem{theorem}{Theorem}[section]
\newtheorem{lemma}[theorem]{Lemma}
\newtheorem{corollary}[theorem]{Corollary}
\newtheorem{proposition}[theorem]{Proposition}

\theoremstyle{definition}
\newtheorem{definition}[theorem]{Definition}
\newtheorem{example}[theorem]{Example}
\newtheorem{examples}[theorem]{Examples}
\newtheorem{remark}[theorem]{Remark}

\numberwithin{equation}{section}

\makeatletter
\newcommand\leanlinksitem[3]{\@namedef{leanlinks@#1}{\leanlinks@box{#2}{#3}}}
\newcommand\leanlinks@box[2]{\ {\normalfont\footnotesize\sffamily[\textcolor{teal}{Lean}#1%
  \if\relax\detokenize{#2}\relax\else; \textit{#2}\fi]\leanlinks@legend}}
\let\leanlinks@legend\relax
\newcommand\leanlinkslegend[3]{\gdef\leanlinks@legend{\footnote{\normalfont\rmfamily
  The boxes [Lean 1 2 \dots] link each numbered item to its Lean declarations in
  \url{#1} at commit \texttt{\detokenize{#2}}. A Definition's declarations are marked
  P (defined in the formalization), M (Mathlib's) or B (defined in the formalization and proved
  equivalent to Mathlib's). A word after the links says how the Lean statement differs from the
  paper's (stronger, weaker, different, or scope); otherwise it is the paper's. At most four
  links are shown; ``+$n$'' leads to the full list in the \href{#3}{index}.}%
  \global\let\leanlinks@legend\relax}}
\newcommand\leanlinkmore[2]{\ \href{#1}{+#2}}
\newcommand\leanlink[3]{\if\relax\detokenize{#2}\relax\href{#1}{\,$\nearrow$}\else
  \ \href{#1}{#2}\fi\if\relax\detokenize{#3}\relax\else\textsuperscript{#3}\fi}
\newcommand\leanlinks@patch{%
  \let\leanlinks@thmhead\thmhead@plain
  \def\thmhead@plain##1##2##3{\leanlinks@thmhead{##1}{##2}{##3}%
    \ifcsname leanlinks@##2\endcsname\csname leanlinks@##2\endcsname\fi}%
  \let\thmhead\thmhead@plain}
\IfFileExists{lean-links.tex}{\leanlinkslegend{https://github.com/DavidSForsman/ulmer-bialgebras-lean}{v1.0.0}{https://github.com/DavidSForsman/ulmer-bialgebras-lean/blob/v1.0.0/PAPER\_LINKS.md}
\leanlinksitem{2.2}{\leanlink{https://github.com/DavidSForsman/ulmer-bialgebras-lean/blob/v1.0.0/UlmerBialgebras/Paper/Sec2.lean\#L68}{1}{}\leanlink{https://github.com/DavidSForsman/ulmer-bialgebras-lean/blob/v1.0.0/UlmerBialgebras/Paper/Sec2.lean\#L75}{2}{}}{partial}
\leanlinksitem{2.3}{\leanlink{https://github.com/DavidSForsman/ulmer-bialgebras-lean/blob/v1.0.0/UlmerBialgebras/Paper/Sec2.lean\#L60}{1}{}\leanlink{https://github.com/DavidSForsman/ulmer-bialgebras-lean/blob/v1.0.0/UlmerBialgebras/Paper/Sec2.lean\#L96}{2}{}\leanlink{https://github.com/DavidSForsman/ulmer-bialgebras-lean/blob/v1.0.0/UlmerBialgebras/Paper/Sec2.lean\#L110}{3}{}\leanlink{https://github.com/DavidSForsman/ulmer-bialgebras-lean/blob/v1.0.0/UlmerBialgebras/Paper/Sec2.lean\#L127}{4}{}\leanlinkmore{https://github.com/DavidSForsman/ulmer-bialgebras-lean/blob/v1.0.0/PAPER\_LINKS.md\#item-2-3}{1}}{partial}
\leanlinksitem{3.2}{\leanlink{https://github.com/DavidSForsman/ulmer-bialgebras-lean/blob/v1.0.0/UlmerBialgebras/Paper/Sec3.lean\#L67}{1}{}\leanlink{https://github.com/DavidSForsman/ulmer-bialgebras-lean/blob/v1.0.0/UlmerBialgebras/Paper/Sec3.lean\#L82}{2}{}}{}
\leanlinksitem{3.11}{\leanlink{https://github.com/DavidSForsman/ulmer-bialgebras-lean/blob/v1.0.0/UlmerBialgebras/Paper/Sec3.lean\#L96}{1}{}\leanlink{https://github.com/DavidSForsman/ulmer-bialgebras-lean/blob/v1.0.0/UlmerBialgebras/Paper/Sec3.lean\#L100}{2}{}\leanlink{https://github.com/DavidSForsman/ulmer-bialgebras-lean/blob/v1.0.0/UlmerBialgebras/Paper/Sec3.lean\#L105}{3}{}}{}
\leanlinksitem{3.12}{\leanlink{https://github.com/DavidSForsman/ulmer-bialgebras-lean/blob/v1.0.0/UlmerBialgebras/Paper/Sec3Ex.lean\#L50}{1}{}\leanlink{https://github.com/DavidSForsman/ulmer-bialgebras-lean/blob/v1.0.0/UlmerBialgebras/Paper/Sec3Ex.lean\#L63}{2}{}\leanlink{https://github.com/DavidSForsman/ulmer-bialgebras-lean/blob/v1.0.0/UlmerBialgebras/Paper/Sec3ExBicat.lean\#L91}{3}{}\leanlink{https://github.com/DavidSForsman/ulmer-bialgebras-lean/blob/v1.0.0/UlmerBialgebras/Paper/Sec3ExBicat.lean\#L143}{4}{}}{stronger}
\leanlinksitem{3.15}{\leanlink{https://github.com/DavidSForsman/ulmer-bialgebras-lean/blob/v1.0.0/UlmerBialgebras/Paper/Sec3.lean\#L110}{}{}}{}
\leanlinksitem{3.17}{\leanlink{https://github.com/DavidSForsman/ulmer-bialgebras-lean/blob/v1.0.0/UlmerBialgebras/Paper/Sec3MonadAlgObject.lean\#L69}{1}{}\leanlink{https://github.com/DavidSForsman/ulmer-bialgebras-lean/blob/v1.0.0/UlmerBialgebras/Paper/Sec3MonadAlgObject.lean\#L95}{2}{}\leanlink{https://github.com/DavidSForsman/ulmer-bialgebras-lean/blob/v1.0.0/UlmerBialgebras/Paper/Sec3MonadAlgObject.lean\#L144}{3}{}}{}
\leanlinksitem{3.18}{\leanlink{https://github.com/DavidSForsman/ulmer-bialgebras-lean/blob/v1.0.0/UlmerBialgebras/Paper/Sec3.lean\#L120}{1}{}\leanlink{https://github.com/DavidSForsman/ulmer-bialgebras-lean/blob/v1.0.0/UlmerBialgebras/Paper/Sec3.lean\#L132}{2}{}}{stronger}
\leanlinksitem{3.20}{\leanlink{https://github.com/DavidSForsman/ulmer-bialgebras-lean/blob/v1.0.0/UlmerBialgebras/Paper/Sec3.lean\#L267}{1}{}\leanlink{https://github.com/DavidSForsman/ulmer-bialgebras-lean/blob/v1.0.0/UlmerBialgebras/Paper/Sec3.lean\#L285}{2}{}\leanlink{https://github.com/DavidSForsman/ulmer-bialgebras-lean/blob/v1.0.0/UlmerBialgebras/Paper/Sec3.lean\#L300}{3}{}\leanlink{https://github.com/DavidSForsman/ulmer-bialgebras-lean/blob/v1.0.0/UlmerBialgebras/Paper/Sec3.lean\#L312}{4}{}\leanlinkmore{https://github.com/DavidSForsman/ulmer-bialgebras-lean/blob/v1.0.0/PAPER\_LINKS.md\#item-3-20}{5}}{stronger}
\leanlinksitem{3.21}{\leanlink{https://github.com/DavidSForsman/ulmer-bialgebras-lean/blob/v1.0.0/UlmerBialgebras/Paper/Sec3.lean\#L143}{1}{}\leanlink{https://github.com/DavidSForsman/ulmer-bialgebras-lean/blob/v1.0.0/UlmerBialgebras/Paper/Sec3.lean\#L155}{2}{}}{}
\leanlinksitem{3.22}{\leanlink{https://github.com/DavidSForsman/ulmer-bialgebras-lean/blob/v1.0.0/UlmerBialgebras/Paper/Sec3.lean\#L195}{}{}}{}
\leanlinksitem{3.23}{\leanlink{https://github.com/DavidSForsman/ulmer-bialgebras-lean/blob/v1.0.0/UlmerBialgebras/Paper/Sec3.lean\#L207}{1}{}\leanlink{https://github.com/DavidSForsman/ulmer-bialgebras-lean/blob/v1.0.0/UlmerBialgebras/Paper/Sec3.lean\#L220}{2}{}\leanlink{https://github.com/DavidSForsman/ulmer-bialgebras-lean/blob/v1.0.0/UlmerBialgebras/Paper/Sec3.lean\#L231}{3}{}}{}
\leanlinksitem{3.24}{\leanlink{https://github.com/DavidSForsman/ulmer-bialgebras-lean/blob/v1.0.0/UlmerBialgebras/Paper/Sec3.lean\#L250}{}{}}{different}
\leanlinksitem{4.5}{\leanlink{https://github.com/DavidSForsman/ulmer-bialgebras-lean/blob/v1.0.0/UlmerBialgebras/Paper/Sec4.lean\#L81}{}{}}{}
\leanlinksitem{4.6}{\leanlink{https://github.com/DavidSForsman/ulmer-bialgebras-lean/blob/v1.0.0/UlmerBialgebras/Paper/Sec4.lean\#L111}{1}{}\leanlink{https://github.com/DavidSForsman/ulmer-bialgebras-lean/blob/v1.0.0/UlmerBialgebras/Paper/Sec4.lean\#L147}{2}{}\leanlink{https://github.com/DavidSForsman/ulmer-bialgebras-lean/blob/v1.0.0/UlmerBialgebras/Paper/Sec4.lean\#L176}{3}{}}{}
\leanlinksitem{4.7}{\leanlink{https://github.com/DavidSForsman/ulmer-bialgebras-lean/blob/v1.0.0/UlmerBialgebras/Paper/Sec4.lean\#L202}{1}{}\leanlink{https://github.com/DavidSForsman/ulmer-bialgebras-lean/blob/v1.0.0/UlmerBialgebras/Paper/Sec4.lean\#L226}{2}{}\leanlink{https://github.com/DavidSForsman/ulmer-bialgebras-lean/blob/v1.0.0/UlmerBialgebras/Paper/Sec4.lean\#L253}{3}{}\leanlink{https://github.com/DavidSForsman/ulmer-bialgebras-lean/blob/v1.0.0/UlmerBialgebras/Paper/Sec4.lean\#L286}{4}{}\leanlinkmore{https://github.com/DavidSForsman/ulmer-bialgebras-lean/blob/v1.0.0/PAPER\_LINKS.md\#item-4-7}{3}}{}
\leanlinksitem{4.8}{\leanlink{https://github.com/DavidSForsman/ulmer-bialgebras-lean/blob/v1.0.0/UlmerBialgebras/Paper/Sec4.lean\#L384}{1}{}\leanlink{https://github.com/DavidSForsman/ulmer-bialgebras-lean/blob/v1.0.0/UlmerBialgebras/Paper/Sec4.lean\#L406}{2}{}}{}
\leanlinksitem{4.9}{\leanlink{https://github.com/DavidSForsman/ulmer-bialgebras-lean/blob/v1.0.0/UlmerBialgebras/Paper/Sec4RelativelyAlgebraic.lean\#L35}{1}{}\leanlink{https://github.com/DavidSForsman/ulmer-bialgebras-lean/blob/v1.0.0/UlmerBialgebras/Paper/Sec4RelativelyAlgebraic.lean\#L40}{2}{}\leanlink{https://github.com/DavidSForsman/ulmer-bialgebras-lean/blob/v1.0.0/UlmerBialgebras/Paper/Sec4RelativelyAlgebraic.lean\#L55}{3}{}}{}
\leanlinksitem{4.10}{\leanlink{https://github.com/DavidSForsman/ulmer-bialgebras-lean/blob/v1.0.0/UlmerBialgebras/Paper/Sec4.lean\#L427}{}{}}{}
\leanlinksitem{4.11}{\leanlink{https://github.com/DavidSForsman/ulmer-bialgebras-lean/blob/v1.0.0/UlmerBialgebras/Paper/Sec4.lean\#L465}{1}{}\leanlink{https://github.com/DavidSForsman/ulmer-bialgebras-lean/blob/v1.0.0/UlmerBialgebras/Paper/Sec4.lean\#L489}{2}{}\leanlink{https://github.com/DavidSForsman/ulmer-bialgebras-lean/blob/v1.0.0/UlmerBialgebras/Paper/Sec4.lean\#L513}{3}{}}{}
\leanlinksitem{4.12}{\leanlink{https://github.com/DavidSForsman/ulmer-bialgebras-lean/blob/v1.0.0/UlmerBialgebras/Paper/Sec4.lean\#L534}{1}{}\leanlink{https://github.com/DavidSForsman/ulmer-bialgebras-lean/blob/v1.0.0/UlmerBialgebras/Paper/Sec4.lean\#L556}{2}{}}{stronger}
\leanlinksitem{4.13}{\leanlink{https://github.com/DavidSForsman/ulmer-bialgebras-lean/blob/v1.0.0/UlmerBialgebras/Paper/Sec4.lean\#L589}{}{}}{}
\leanlinksitem{4.14}{\leanlink{https://github.com/DavidSForsman/ulmer-bialgebras-lean/blob/v1.0.0/UlmerBialgebras/Paper/Sec4.lean\#L620}{}{}}{}
\leanlinksitem{4.15}{\leanlink{https://github.com/DavidSForsman/ulmer-bialgebras-lean/blob/v1.0.0/UlmerBialgebras/Paper/Sec4.lean\#L653}{}{}}{}
\leanlinksitem{4.16}{\leanlink{https://github.com/DavidSForsman/ulmer-bialgebras-lean/blob/v1.0.0/UlmerBialgebras/Paper/Sec4.lean\#L684}{}{}}{}
\leanlinksitem{4.17}{\leanlink{https://github.com/DavidSForsman/ulmer-bialgebras-lean/blob/v1.0.0/UlmerBialgebras/Paper/Sec4.lean\#L707}{}{}}{}
\leanlinksitem{4.18}{\leanlink{https://github.com/DavidSForsman/ulmer-bialgebras-lean/blob/v1.0.0/UlmerBialgebras/Paper/Sec4.lean\#L734}{}{}}{}
\leanlinksitem{4.19}{\leanlink{https://github.com/DavidSForsman/ulmer-bialgebras-lean/blob/v1.0.0/UlmerBialgebras/Paper/Sec4.lean\#L761}{}{}}{}
\leanlinksitem{4.20}{\leanlink{https://github.com/DavidSForsman/ulmer-bialgebras-lean/blob/v1.0.0/UlmerBialgebras/Paper/Sec4Hopf.lean\#L74}{1}{}\leanlink{https://github.com/DavidSForsman/ulmer-bialgebras-lean/blob/v1.0.0/UlmerBialgebras/Paper/Sec4Hopf.lean\#L106}{2}{}\leanlink{https://github.com/DavidSForsman/ulmer-bialgebras-lean/blob/v1.0.0/UlmerBialgebras/Paper/Sec4Hopf.lean\#L143}{3}{}\leanlink{https://github.com/DavidSForsman/ulmer-bialgebras-lean/blob/v1.0.0/UlmerBialgebras/Paper/Sec4Hopf.lean\#L183}{4}{}\leanlinkmore{https://github.com/DavidSForsman/ulmer-bialgebras-lean/blob/v1.0.0/PAPER\_LINKS.md\#item-4-20}{5}}{stronger, different}
\leanlinksitem{4.21}{\leanlink{https://github.com/DavidSForsman/ulmer-bialgebras-lean/blob/v1.0.0/UlmerBialgebras/Paper/Sec4.lean\#L795}{}{}}{}
\leanlinksitem{5.2}{\leanlink{https://github.com/DavidSForsman/ulmer-bialgebras-lean/blob/v1.0.0/UlmerBialgebras/Paper/Sec5.lean\#L106}{}{}}{}
\leanlinksitem{5.3}{\leanlink{https://github.com/DavidSForsman/ulmer-bialgebras-lean/blob/v1.0.0/UlmerBialgebras/Paper/Sec5.lean\#L139}{1}{}\leanlink{https://github.com/DavidSForsman/ulmer-bialgebras-lean/blob/v1.0.0/UlmerBialgebras/Paper/Sec5.lean\#L175}{2}{}}{stronger}
\leanlinksitem{5.5}{\leanlink{https://github.com/DavidSForsman/ulmer-bialgebras-lean/blob/v1.0.0/UlmerBialgebras/Paper/Sec5.lean\#L210}{1}{}\leanlink{https://github.com/DavidSForsman/ulmer-bialgebras-lean/blob/v1.0.0/UlmerBialgebras/Paper/Sec5.lean\#L240}{2}{}}{}
\leanlinksitem{5.6}{\leanlink{https://github.com/DavidSForsman/ulmer-bialgebras-lean/blob/v1.0.0/UlmerBialgebras/Paper/Sec5.lean\#L286}{1}{}\leanlink{https://github.com/DavidSForsman/ulmer-bialgebras-lean/blob/v1.0.0/UlmerBialgebras/Paper/Sec5.lean\#L328}{2}{}}{}
\leanlinksitem{5.7}{\leanlink{https://github.com/DavidSForsman/ulmer-bialgebras-lean/blob/v1.0.0/UlmerBialgebras/Paper/Sec5.lean\#L378}{}{}}{different}
\leanlinksitem{5.8}{\leanlink{https://github.com/DavidSForsman/ulmer-bialgebras-lean/blob/v1.0.0/UlmerBialgebras/Paper/Sec5.lean\#L420}{}{}}{different}
\leanlinksitem{5.9}{\leanlink{https://github.com/DavidSForsman/ulmer-bialgebras-lean/blob/v1.0.0/UlmerBialgebras/Paper/Sec5UniversalAlgebra.lean\#L76}{1}{}\leanlink{https://github.com/DavidSForsman/ulmer-bialgebras-lean/blob/v1.0.0/UlmerBialgebras/Paper/Sec5UniversalAlgebra.lean\#L101}{2}{}\leanlink{https://github.com/DavidSForsman/ulmer-bialgebras-lean/blob/v1.0.0/UlmerBialgebras/Paper/Sec5.lean\#L468}{3}{}}{stronger, different}
\leanlinksitem{5.10}{\leanlink{https://github.com/DavidSForsman/ulmer-bialgebras-lean/blob/v1.0.0/UlmerBialgebras/Paper/Sec5.lean\#L502}{1}{}\leanlink{https://github.com/DavidSForsman/ulmer-bialgebras-lean/blob/v1.0.0/UlmerBialgebras/Paper/Sec5.lean\#L527}{2}{}\leanlink{https://github.com/DavidSForsman/ulmer-bialgebras-lean/blob/v1.0.0/UlmerBialgebras/Paper/Sec5.lean\#L606}{3}{}}{stronger}
\leanlinksitem{5.11}{\leanlink{https://github.com/DavidSForsman/ulmer-bialgebras-lean/blob/v1.0.0/UlmerBialgebras/Paper/Sec5.lean\#L562}{1}{}\leanlink{https://github.com/DavidSForsman/ulmer-bialgebras-lean/blob/v1.0.0/UlmerBialgebras/Paper/Sec5.lean\#L638}{2}{}\leanlink{https://github.com/DavidSForsman/ulmer-bialgebras-lean/blob/v1.0.0/UlmerBialgebras/Paper/Sec5.lean\#L672}{3}{}}{stronger}
\leanlinksitem{6.2}{\leanlink{https://github.com/DavidSForsman/ulmer-bialgebras-lean/blob/v1.0.0/UlmerBialgebras/Paper/Sec6.lean\#L103}{}{}}{}
\leanlinksitem{6.3}{\leanlink{https://github.com/DavidSForsman/ulmer-bialgebras-lean/blob/v1.0.0/UlmerBialgebras/Paper/Sec6.lean\#L130}{1}{}\leanlink{https://github.com/DavidSForsman/ulmer-bialgebras-lean/blob/v1.0.0/UlmerBialgebras/Paper/Sec6.lean\#L160}{2}{}\leanlink{https://github.com/DavidSForsman/ulmer-bialgebras-lean/blob/v1.0.0/UlmerBialgebras/Paper/Sec6.lean\#L191}{3}{}}{}
\leanlinksitem{6.4}{\leanlink{https://github.com/DavidSForsman/ulmer-bialgebras-lean/blob/v1.0.0/UlmerBialgebras/Paper/Sec6.lean\#L216}{1}{}\leanlink{https://github.com/DavidSForsman/ulmer-bialgebras-lean/blob/v1.0.0/UlmerBialgebras/Paper/Sec6.lean\#L243}{2}{}}{}
\leanlinksitem{6.5}{\leanlink{https://github.com/DavidSForsman/ulmer-bialgebras-lean/blob/v1.0.0/UlmerBialgebras/Paper/Sec6.lean\#L272}{}{}}{}
\leanlinksitem{6.6}{\leanlink{https://github.com/DavidSForsman/ulmer-bialgebras-lean/blob/v1.0.0/UlmerBialgebras/Paper/Sec6.lean\#L313}{1}{}\leanlink{https://github.com/DavidSForsman/ulmer-bialgebras-lean/blob/v1.0.0/UlmerBialgebras/Paper/Sec6.lean\#L359}{2}{}\leanlink{https://github.com/DavidSForsman/ulmer-bialgebras-lean/blob/v1.0.0/UlmerBialgebras/Paper/Sec6.lean\#L400}{3}{}\leanlink{https://github.com/DavidSForsman/ulmer-bialgebras-lean/blob/v1.0.0/UlmerBialgebras/Paper/Sec6.lean\#L435}{4}{}}{stronger, different}
\leanlinksitem{6.7}{\leanlink{https://github.com/DavidSForsman/ulmer-bialgebras-lean/blob/v1.0.0/UlmerBialgebras/Paper/Sec6.lean\#L470}{}{}}{different}
\leanlinksitem{6.8}{\leanlink{https://github.com/DavidSForsman/ulmer-bialgebras-lean/blob/v1.0.0/UlmerBialgebras/Paper/Sec6.lean\#L518}{}{}}{different}
\leanlinksitem{6.9}{\leanlink{https://github.com/DavidSForsman/ulmer-bialgebras-lean/blob/v1.0.0/UlmerBialgebras/Paper/Sec6.lean\#L565}{}{}}{stronger, different}
\leanlinksitem{6.10}{\leanlink{https://github.com/DavidSForsman/ulmer-bialgebras-lean/blob/v1.0.0/UlmerBialgebras/Paper/Sec6Loc.lean\#L72}{}{}}{stronger}
\leanlinksitem{6.11}{\leanlink{https://github.com/DavidSForsman/ulmer-bialgebras-lean/blob/v1.0.0/UlmerBialgebras/Paper/Sec6Hopf.lean\#L95}{1}{}\leanlink{https://github.com/DavidSForsman/ulmer-bialgebras-lean/blob/v1.0.0/UlmerBialgebras/Paper/Sec6Hopf.lean\#L150}{2}{}\leanlink{https://github.com/DavidSForsman/ulmer-bialgebras-lean/blob/v1.0.0/UlmerBialgebras/Paper/Sec6Hopf.lean\#L185}{3}{}}{stronger, different}
\leanlinksitem{7.1}{\leanlink{https://github.com/DavidSForsman/ulmer-bialgebras-lean/blob/v1.0.0/UlmerBialgebras/Paper/Sec7.lean\#L97}{1}{}\leanlink{https://github.com/DavidSForsman/ulmer-bialgebras-lean/blob/v1.0.0/UlmerBialgebras/Paper/Sec7Variant.lean\#L62}{2}{}}{}
\leanlinksitem{7.4}{\leanlink{https://github.com/DavidSForsman/ulmer-bialgebras-lean/blob/v1.0.0/UlmerBialgebras/Paper/Sec7.lean\#L117}{}{}}{stronger}
\leanlinksitem{7.5}{\leanlink{https://github.com/DavidSForsman/ulmer-bialgebras-lean/blob/v1.0.0/UlmerBialgebras/Paper/Sec7.lean\#L142}{}{}}{stronger}
\leanlinksitem{7.6}{\leanlink{https://github.com/DavidSForsman/ulmer-bialgebras-lean/blob/v1.0.0/UlmerBialgebras/Paper/Sec7.lean\#L182}{1}{}\leanlink{https://github.com/DavidSForsman/ulmer-bialgebras-lean/blob/v1.0.0/UlmerBialgebras/Paper/Sec7.lean\#L225}{2}{}}{stronger}
\leanlinksitem{7.7}{\leanlink{https://github.com/DavidSForsman/ulmer-bialgebras-lean/blob/v1.0.0/UlmerBialgebras/Paper/Sec7.lean\#L264}{1}{}\leanlink{https://github.com/DavidSForsman/ulmer-bialgebras-lean/blob/v1.0.0/UlmerBialgebras/Paper/Sec7.lean\#L301}{2}{}}{stronger, different}
\leanlinksitem{A.3}{\leanlink{https://github.com/DavidSForsman/ulmer-bialgebras-lean/blob/v1.0.0/UlmerBialgebras/Paper/Appendix.lean\#L104}{1}{}\leanlink{https://github.com/DavidSForsman/ulmer-bialgebras-lean/blob/v1.0.0/UlmerBialgebras/Paper/AppendixA3A5.lean\#L73}{2}{}\leanlink{https://github.com/DavidSForsman/ulmer-bialgebras-lean/blob/v1.0.0/UlmerBialgebras/Paper/AppendixA3A5.lean\#L121}{3}{}}{}
\leanlinksitem{A.4}{\leanlink{https://github.com/DavidSForsman/ulmer-bialgebras-lean/blob/v1.0.0/UlmerBialgebras/Paper/Appendix.lean\#L141}{}{}}{stronger}
\leanlinksitem{A.5}{\leanlink{https://github.com/DavidSForsman/ulmer-bialgebras-lean/blob/v1.0.0/UlmerBialgebras/Paper/Appendix.lean\#L68}{1}{}\leanlink{https://github.com/DavidSForsman/ulmer-bialgebras-lean/blob/v1.0.0/UlmerBialgebras/Paper/AppendixA3A5.lean\#L100}{2}{}}{different}
\leanlinks@patch}{}
\makeatother
